\documentclass[11pt,letterpaper]{article}

\usepackage[margin=0.6in]{geometry}
\usepackage{enumitem}
\usepackage{titlesec}
\usepackage{hyperref}
\usepackage[T1]{fontenc}
\usepackage[dvipsnames]{xcolor}
\usepackage{tabularx}
\usepackage{mathpazo}       % Palatino — a timeless serif typeface
\usepackage[stretch=10,shrink=10]{microtype}  % Refined character spacing

% --- Colour palette (traditional dark navy) ---
\definecolor{accentnavy}{RGB}{10, 40, 80}
\definecolor{textdark}{RGB}{30, 30, 30}

\hypersetup{
    colorlinks=true,
    urlcolor=accentnavy,
    linkcolor=accentnavy
}

% --- Section headers: bold small-caps with rule ---
\titleformat{\section}
{\large\bfseries\scshape\color{accentnavy}}
{}
{0em}
{}
[\titlerule]

\titlespacing{\section}{0pt}{10pt}{5pt}

\pagestyle{empty}
\setlength{\parindent}{0pt}

\setlist[itemize]{
leftmargin=1.3em,
itemsep=1pt,
topsep=2pt
}

\begin{document}

%================================================
% HEADER
%================================================

\begin{center}

{\LARGE\bfseries\scshape Venkata Sai Manideep Cherukuri}

\vspace{4pt}


\vspace{6pt}

{\small
\href{mailto:saimanideep.ch12345@gmail.com}{saimanideep.ch12345@gmail.com}
\quad\textbar\quad
+91\,80963\,68241
\quad\textbar\quad
\href{https://github.com/manideep1428}{github.com/manideep1428}
}

\end{center}

%================================================
% EXPERIENCE
%================================================

\section{Experience}

\textbf{Emergx}
\hfill
\textit{May 2025 -- Sept 2025}

Software Development Engineer Intern \hfill Remote

{\small \href{https://www.emergx.ai}{emergx.ai}}

\begin{itemize}
    \item Built and shipped a full-stack AI platform supporting supply-chain workflows and automation.
    \item Developed frontend and backend systems used in production environments.
    \item Improved application scalability and reliability as the platform grew to over 5,000 users.
\end{itemize}

\vspace{4pt}

\textbf{Freelance Technical Writer}
\hfill
\textit{2024 -- 2025}

Mintlify Documentation

\begin{itemize}
    \item Authored developer documentation for Brew, Emby, and Debounce.
    \item Wrote API references, onboarding guides, authentication flows, and integration tutorials.
    \item Improved developer onboarding and product adoption through structured documentation.
\end{itemize}

%================================================
% EDUCATION
%================================================

\section{Education}

\textbf{Loyola College}
\hfill
\textit{2024 -- 2026}

Bachelor of Computer Science (Dropout) \hfill Guntur, India

%================================================
% PROJECTS
%================================================

\section{Projects}

\textbf{Kiver}
\hfill
\textit{Feb 2026 -- Present}

Desktop Mission Control for AI Coding Agents

{\small
\href{https://github.com/manideep1428/Kiver}{GitHub}
}

\begin{itemize}
    \item Built a unified control plane and visual workspace for local AI coding agents (Claude Code and Codex CLI).
    \item Coordinated parallel autonomous subagents with deterministic SQLite file locking to prevent concurrent file conflicts.
    \item Implemented live terminal multiplexing with xterm.js/node-pty, zero-loss mid-task provider migration, and per-agent git rollbacks.
\end{itemize}

\vspace{5pt}

\textbf{Supergent}
\hfill
\textit{Mar 2026 -- Present}

AI Full-Stack App Builder

{\small
\href{https://supergent-web.vercel.app/}{Live Demo}
\quad\textbar\quad
\href{https://github.com/manideep1428/supergent}{GitHub}
}

\begin{itemize}
    \item Built an AI software engineer capable of generating complete web and mobile applications from a single prompt.
    \item Generates frontend, backend APIs, authentication, databases, and deployment configurations automatically.
    \item Implemented multi-agent planning, coding, debugging, and refinement workflows.
\end{itemize}

\vspace{5pt}

\textbf{LovDacn UI}
\hfill
\textit{2026 -- Present}

Open Source UI Component Library (Beta)

{\small
\href{https://github.com/lovdacn-ui/ui}{github.com/lovdacn-ui/ui}
}

\begin{itemize}
    \item Creator and maintainer of an open-source React and Next.js UI component library.
    \item Building reusable UI components, templates, and developer tools using TypeScript and TailwindCSS.
    \item Currently in beta and actively expanding the component collection and documentation.
\end{itemize}

\vspace{5pt}

\textbf{Instagram MCP}
\hfill
\textit{2026 -- Present}

MCP Server for Instagram Automation

{\small
\href{https://github.com/manideep1428/instagram_MCP}{GitHub}
}

\begin{itemize}
    \item Built a Model Context Protocol (MCP) server that lets AI agents automate Instagram through natural language commands.
    \item Integrated the Instagram Graph API with MCP-compatible clients such as Claude Desktop, Cursor, and Windsurf.
    \item Implemented tools for media publishing, profile and comment management, direct messaging, and account automation through a unified interface.
\end{itemize}

\vspace{5pt}

\textbf{Niana}
\hfill
\textit{Jan 2026 -- Present}

AI Design Generation Platform

{\small
\href{https://niana.design}{niana.design}
}

\begin{itemize}
    \item Built a platform that generates mobile UI designs from natural language prompts.
    \item Added direct export workflows for Figma integration.
    \item Launched subscription billing and onboarded paying users.
\end{itemize}

\vspace{5pt}

\textbf{Astron}
\hfill
\textit{Feb 2026 -- Present}

AI Browser Automation CLI

{\small
\href{https://github.com/manideep1428/astron-browser}{GitHub}
}

\begin{itemize}
    \item Built a terminal-native browser automation platform powered by AI agents and natural language commands.
    \item Integrated OpenAI, Claude, Gemini, and Browser Use with seamless model switching.
    \item Added persistent daemon mode, browser session management, and secure credential storage.
\end{itemize}

\vspace{5pt}


\textbf{Auto AI Video Generator}
\hfill
\textit{Dec 2024 -- Jan 2025}

{\small
\href{https://github.com/manideep1428/auto_ai_video_generator}{GitHub}
}

\begin{itemize}
    \item Built an automated AI pipeline that creates videos from a single topic.
    \item Automated script writing, image generation, voice synthesis, editing, and publishing.
    \item Generated over 140,000+ views during the first week after launch.
\end{itemize}

\vspace{5pt}

\textbf{Cheat-Coder}
\hfill
\textit{Dec 2024 -- Jul 2025}

AI Interview Assistant

{\small
\href{https://github.com/manideep1428/cheat-coder}{GitHub}
}

\begin{itemize}
    \item Built a local-first AI interview assistant powered by user-provided API keys.
    \item Designed a privacy-focused architecture with zero server-side storage.
    \item Tested across hundreds of users in real interview environments.
\end{itemize}

\section{Achievements}
\vspace{6pt}

\begin{itemize}
    \item Participated and earned certification in \textbf{Mumbai Hack 2025 Hackathon}, demonstrating strong problem-solving and teamwork skills.
\end{itemize}
\vspace{6pt}
%================================================
% SKILLS
%================================================

\section{Skills}

\renewcommand{\arraystretch}{1.25}

\begin{tabularx}{\linewidth}{@{}lX@{}}

\textbf{Programming} &
Python, TypeScript, JavaScript, SQL, HTML, CSS \\

\textbf{Frontend} &
React, Next.js, React Native, TailwindCSS \\

\textbf{Backend} &
Node.js, PostgreSQL, MongoDB, Prisma ORM, REST APIs \\

\textbf{AI / LLMs} &
OpenAI SDK, Anthropic Claude SDK, Google Gemini SDK, Vercel AI SDK, LangChain, LangGraph, Hugging Face Transformers, PyTorch \\

\textbf{Cloud \& DevOps} &
AWS, GCP, Azure, Docker, Kubernetes, Supabase, Convex, Vercel \\

\textbf{Tools} &
Git, GitHub, Linux, Postman, Razorpay \\

\end{tabularx}

%================================================
% PUBLICATIONS
%================================================

\section{Publications}

\begin{itemize}
    \item \href{https://medium.com/@saimanideep.ch12345/why-floats-are-preferred-over-integers-in-machine-learning-73da30ca2225}{Why Floats are Preferred Over Integers in Machine Learning}

    \item \href{https://medium.com/@saimanideep.ch12345/genesis-training-robots-430-000x-faster-than-real-time-c6b6775eee63}{Genesis: Training Robots 430,000x Faster Than Real-Time}
\end{itemize}

\end{document}